
\documentclass[aspectratio=169]{beamer}
\usetheme{Madrid}
\usecolortheme{default}
\usefonttheme{professionalfonts}

\usepackage[utf8]{inputenc}
\usepackage[T1]{fontenc}
\usepackage{amsmath, amssymb, amsfonts}
\usepackage{booktabs}
\usepackage{multicol}
\usepackage{tikz}
\usepackage{graphicx}
\usepackage{array}
\usepackage{hyperref}

\title[GED Math: Data, Statistics \& Probability]{GED Math Unit:\\Data, Statistics, and Probability}

\institute{Instructor: Paul Fox}
\date{\today}

\hypersetup{
  colorlinks=true,
  urlcolor=blue,
  linkcolor=blue
}

\newcommand{\exampletitle}[1]{\vspace{0.25em}\textbf{Example.} #1\vspace{0.25em}}
\newcommand{\worked}{\textbf{Worked Solution.}}

\begin{document}

% Title
\begin{frame}
  \titlepage
\end{frame}

% Table of Contents
\begin{frame}{Unit Map}
  \tableofcontents
\end{frame}

%%%%%%%%%%%%%%%%%%%%%%%%%%%%%%%%%%%%%%%%%%%%%%%%%%%%%%%%%%%%
\section{Lesson 1: Data Types, Tables, and Sampling}
%%%%%%%%%%%%%%%%%%%%%%%%%%%%%%%%%%%%%%%%%%%%%%%%%%%%%%%%%%%%

\begin{frame}{Objectives}
\begin{itemize}
  \item Identify data types: categorical vs. numerical (discrete/continuous).
  \item Organize data using frequency and relative-frequency tables.
  \item Understand populations, samples, and basic sampling methods.
  \item Recognize sources of bias and good survey practices.
\end{itemize}
\end{frame}

\begin{frame}{Key Ideas}
\begin{itemize}
  \item \textbf{Categorical} data: labels or groups (e.g., color, brand).
  \item \textbf{Numerical} data: numbers that measure/count.
    \begin{itemize}
      \item \textbf{Discrete}: counts (0,1,2,\ldots).
      \item \textbf{Continuous}: measured on a scale (height, time).
    \end{itemize}
  \item \textbf{Frequency table}: lists categories/values with counts.
  \item \textbf{Relative frequency}: frequency divided by total.
  \item \textbf{Population}: entire group; \textbf{sample}: subset.
\end{itemize}
\end{frame}

\begin{frame}{Sampling Methods \& Bias}
\begin{itemize}
  \item \textbf{Simple random sample}: every member equally likely.
  \item \textbf{Systematic}: every $k$th individual.
  \item \textbf{Stratified}: split into subgroups, sample from each.
  \item \textbf{Cluster}: randomly select groups, survey all within.
  \item \textbf{Convenience/Voluntary response}: often biased.
  \item \textbf{Bias} occurs when sample or questions systematically favor an outcome.
\end{itemize}
\end{frame}

% Worked Examples (5)
\begin{frame}{Worked Example 1}
\exampletitle{Classify each: (a) brand of phone; (b) number of siblings; (c) commute time.}\\[0.25em]
\worked
\begin{enumerate}
  \item (a) \textbf{Categorical} (brand is a label).
  \item (b) \textbf{Numerical discrete} (counts: 0,1,2,\dots).
  \item (c) \textbf{Numerical continuous} (time measured on a scale).
\end{enumerate}
\end{frame}

\begin{frame}{Worked Example 2}
\exampletitle{Create a frequency and relative-frequency table for colors: R, B, B, G, B, R, G (R=Red,B=Blue,G=Green).}\\[0.25em]
\worked
\begin{enumerate}
  \item Count: R:2,\; B:3,\; G:2 (total $n=7$).
  \item Relative freq: R:$2/7\approx0.286$, B:$3/7\approx0.429$, G:$2/7\approx0.286$.
\end{enumerate}
\end{frame}

\begin{frame}{Worked Example 3}
\exampletitle{A school of 1200 students wants average screen time. Suggest a sampling plan that limits bias.}\\[0.25em]
\worked
\begin{enumerate}
  \item Use a \textbf{stratified random sample} by grade: 9--12.
  \item From each grade, randomly select proportional numbers (e.g., 25 from each if grades are similar size).
  \item Ensure \textbf{neutral wording} and anonymous responses to reduce response bias.
\end{enumerate}
\end{frame}

\begin{frame}{Worked Example 4}
\exampletitle{In a production line, survey every 10th item. What method is this and a potential risk?}\\[0.25em]
\worked
\begin{enumerate}
  \item \textbf{Systematic sampling} with $k=10$.
  \item Risk: if defects occur periodically every 10 items, bias may occur.
\end{enumerate}
\end{frame}

\begin{frame}{Worked Example 5}
\exampletitle{A website poll asks, “Do you agree that students spend too much time on phones?” Identify the issue.}\\[0.25em]
\worked
\begin{enumerate}
  \item \textbf{Voluntary response} (self-selection) and \textbf{leading question} (“too much”) introduce bias.
\end{enumerate}
\end{frame}

%%%%%%%%%%%%%%%%%%%%%%%%%%%%%%%%%%%%%%%%%%%%%%%%%%%%%%%%%%%%
\section{Lesson 2: Graphical Displays (Bar, Histogram, Dot, Pie)}
%%%%%%%%%%%%%%%%%%%%%%%%%%%%%%%%%%%%%%%%%%%%%%%%%%%%%%%%%%%%

\begin{frame}{Objectives}
\begin{itemize}
  \item Construct and interpret bar charts (categorical) and histograms (numerical).
  \item Read dot plots and pie charts; compare distributions.
  \item Describe shape: symmetric, skewed, unimodal/multimodal, gaps, outliers.
\end{itemize}
\end{frame}

\begin{frame}{Key Ideas}
\begin{itemize}
  \item \textbf{Bar chart}: separated bars for categories.
  \item \textbf{Histogram}: contiguous bars for bins of numerical data.
  \item \textbf{Dot plot}: shows each observation as a dot over a number line.
  \item \textbf{Pie chart}: sector angle $\propto$ category fraction.
  \item Look for \textbf{center, spread, shape, outliers}.
\end{itemize}
\end{frame}

% Worked Examples (5)
\begin{frame}{Worked Example 1}
\exampletitle{Categorical data of car colors: Black 18, White 12, Gray 10, Red 5. What display?}\\[0.25em]
\worked
\begin{enumerate}
  \item Use a \textbf{bar chart} (categories).
  \item Heights proportional to counts: 18, 12, 10, 5.
\end{enumerate}
\end{frame}

\begin{frame}{Worked Example 2}
\exampletitle{Create histogram bins for commute times (min) 5,7,9,12,14,16,20,21,22 with bin width 5.}\\[0.25em]
\worked
\begin{enumerate}
  \item Bins: [5,10):3 (5,7,9); [10,15):2 (12,14); [15,20):1 (16); [20,25):3 (20,21,22).
  \item Bars touch because continuous scale.
\end{enumerate}
\end{frame}

\begin{frame}{Worked Example 3}
\exampletitle{A dot plot clusters around 10--12 with a few at 3 and 20. Describe shape.}\\[0.25em]
\worked
\begin{enumerate}
  \item \textbf{Unimodal} center near 11; potential \textbf{outliers} at 3 and 20; mild \textbf{skewness} depending on tail lengths.
\end{enumerate}
\end{frame}

\begin{frame}{Worked Example 4}
\exampletitle{Pie chart sectors for categories A:30, B:45, C:25. Compute angles.}\\[0.25em]
\worked
\begin{enumerate}
  \item Total $=100$. Angles $= \text{fraction}\times 360^\circ$.
  \item A: $0.30\cdot 360^\circ=108^\circ$, B: $162^\circ$, C: $90^\circ$.
\end{enumerate}
\end{frame}

\begin{frame}{Worked Example 5}
\exampletitle{Two histograms: Set X tightly around 50; Set Y spread 30--70. Who has greater variability?}\\[0.25em]
\worked
\begin{enumerate}
  \item Set Y shows greater spread; Set X is more consistent.
\end{enumerate}
\end{frame}

%%%%%%%%%%%%%%%%%%%%%%%%%%%%%%%%%%%%%%%%%%%%%%%%%%%%%%%%%%%%
\section{Lesson 3: Measures of Center and Spread}
%%%%%%%%%%%%%%%%%%%%%%%%%%%%%%%%%%%%%%%%%%%%%%%%%%%%%%%%%%%%

\begin{frame}{Objectives}
\begin{itemize}
  \item Compute mean, median, mode, and range.
  \item Find quartiles, interquartile range (IQR), and identify outliers.
  \item Understand mean absolute deviation (MAD) conceptually.
\end{itemize}
\end{frame}

\begin{frame}{Key Formulas}
\begin{itemize}
  \item Mean: $\displaystyle \bar{x}=\frac{\sum x_i}{n}$.
  \item Median: middle after sort (or average of two middle values).
  \item IQR: $Q_3-Q_1$. \quad Outlier rule: below $Q_1-1.5\text{IQR}$ or above $Q_3+1.5\text{IQR}$.
  \item Range: $\max - \min$. \quad MAD: average of $|x_i-\bar{x}|$.
\end{itemize}
\end{frame}

% Worked Examples (5)
\begin{frame}{Worked Example 1}
\exampletitle{Data: 4, 7, 8, 9, 10. Find mean, median, mode, range.}\\[0.25em]
\worked
\begin{enumerate}
  \item Mean: $(4+7+8+9+10)/5=38/5=7.6$.
  \item Median: middle is 8.
  \item Mode: none (all unique).
  \item Range: $10-4=6$.
\end{enumerate}
\end{frame}

\begin{frame}{Worked Example 2}
\exampletitle{Data: 1, 2, 2, 5, 7, 9. Find median.}\\[0.25em]
\worked
\begin{enumerate}
  \item Even $n=6$: middle two are 2 and 5. Median $=(2+5)/2=3.5$.
\end{enumerate}
\end{frame}

\begin{frame}{Worked Example 3}
\exampletitle{Find $Q_1$, $Q_3$, and IQR for data: 3,4,5,7,8,10,12,13,15.}\\[0.25em]
\worked
\begin{enumerate}
  \item Sorted already, $n=9$. Median is 8 (5th value).
  \item Lower half: 3,4,5,7 $\Rightarrow Q_1$ median $=(4+5)/2=4.5$.
  \item Upper half: 10,12,13,15 $\Rightarrow Q_3$ median $=(12+13)/2=12.5$.
  \item IQR $=12.5-4.5=8$.
\end{enumerate}
\end{frame}

\begin{frame}{Worked Example 4}
\exampletitle{Using the IQR rule, are there outliers in data: 2,3,4,6,7,7,8,50?}\\[0.25em]
\worked
\begin{enumerate}
  \item Sorted: 2,3,4,6,7,7,8,50; $n=8$.
  \item Median $=(6+7)/2=6.5$. Lower: 2,3,4,6 $\Rightarrow Q_1=(3+4)/2=3.5$.
  \item Upper: 7,7,8,50 $\Rightarrow Q_3=(7+8)/2=7.5$.
  \item IQR $=7.5-3.5=4$. Fences: $3.5-1.5(4)=-2.5$ and $7.5+1.5(4)=13.5$.
  \item $50>13.5$ $\Rightarrow$ \textbf{50 is an outlier}.
\end{enumerate}
\end{frame}

\begin{frame}{Worked Example 5}
\exampletitle{Compute the mean and MAD for data: 2, 4, 6, 8.}\\[0.25em]
\worked
\begin{enumerate}
  \item Mean $\bar{x}=(2+4+6+8)/4=5$.
  \item Deviations: $|2-5|=3, |4-5|=1, |6-5|=1, |8-5|=3$.
  \item MAD $=(3+1+1+3)/4=2$.
\end{enumerate}
\end{frame}

%%%%%%%%%%%%%%%%%%%%%%%%%%%%%%%%%%%%%%%%%%%%%%%%%%%%%%%%%%%%
\section{Lesson 4: Boxplots and the Five-Number Summary}
%%%%%%%%%%%%%%%%%%%%%%%%%%%%%%%%%%%%%%%%%%%%%%%%%%%%%%%%%%%%

\begin{frame}{Objectives}
\begin{itemize}
  \item Find the five-number summary: $\min, Q_1, \text{Median}, Q_3, \max$.
  \item Draw and interpret a box-and-whisker plot.
  \item Compare two boxplots to discuss center and spread.
\end{itemize}
\end{frame}

\begin{frame}{Key Ideas}
\begin{itemize}
  \item Box spans from $Q_1$ to $Q_3$; line at the median.
  \item Whiskers extend to min/max that are not outliers.
  \item Outliers plotted as individual points.
  \item IQR measures the middle 50\% spread.
\end{itemize}
\end{frame}

% Worked Examples (5)
\begin{frame}{Worked Example 1}
\exampletitle{Data: 1,3,4,4,6,7,9,11. Find five-number summary.}\\[0.25em]
\worked
\begin{enumerate}
  \item $\min=1$, $\max=11$, $n=8$.
  \item Median $=(4+6)/2=5$.
  \item Lower half: 1,3,4,4 $\Rightarrow Q_1=(3+4)/2=3.5$.
  \item Upper half: 6,7,9,11 $\Rightarrow Q_3=(7+9)/2=8$.
\end{enumerate}
\end{frame}

\begin{frame}{Worked Example 2}
\exampletitle{Sketch the boxplot using the summary from Example 1.}\\[0.25em]
\worked
\begin{enumerate}
  \item Draw number line from 1 to 11.
  \item Box from $Q_1=3.5$ to $Q_3=8$ with median at 5.
  \item Whiskers to 1 and 11 (no outliers by IQR rule checked below).
\end{enumerate}
\end{frame}

\begin{frame}{Worked Example 3}
\exampletitle{Check for outliers using IQR in Example 1.}\\[0.25em]
\worked
\begin{enumerate}
  \item IQR $=8-3.5=4.5$. Fences: $3.5-1.5(4.5)=-3.25$, $8+1.5(4.5)=14.75$.
  \item All data between 1 and 11, so none are outliers.
\end{enumerate}
\end{frame}

\begin{frame}{Worked Example 4}
\exampletitle{Compare two boxplots: Class A median 75, IQR 10; Class B median 72, IQR 20. Who is more consistent?}\\[0.25em]
\worked
\begin{enumerate}
  \item Class A has smaller IQR (10 vs 20) $\Rightarrow$ more consistent (less spread).
  \item Class A also has slightly higher center.
\end{enumerate}
\end{frame}

\begin{frame}{Worked Example 5}
\exampletitle{Given a boxplot with lower whisker at 20, $Q_1=30$, median 40, $Q_3=50$, upper whisker at 65, estimate the percentage of data between 30 and 50.}\\[0.25em]
\worked
\begin{enumerate}
  \item The box from $Q_1$ to $Q_3$ contains the middle 50\% of data.
  \item Therefore, about \textbf{50\%} of data lie between 30 and 50.
\end{enumerate}
\end{frame}

%%%%%%%%%%%%%%%%%%%%%%%%%%%%%%%%%%%%%%%%%%%%%%%%%%%%%%%%%%%%
\section{Lesson 5: Probability Basics}
%%%%%%%%%%%%%%%%%%%%%%%%%%%%%%%%%%%%%%%%%%%%%%%%%%%%%%%%%%%%

\begin{frame}{Objectives}
\begin{itemize}
  \item Compute simple probabilities from equally likely outcomes.
  \item Understand complements, independent vs. dependent events.
  \item Use formulas for compound probability (AND/OR).
\end{itemize}
\end{frame}

\begin{frame}{Key Formulas}
\begin{itemize}
  \item $P(\text{event})=\dfrac{\text{favorable}}{\text{total}}$ (equally likely).
  \item Complement: $P(A^c)=1-P(A)$.
  \item Independent: $P(A\cap B)=P(A)P(B)$.
  \item General OR: $P(A\cup B)=P(A)+P(B)-P(A\cap B)$.
\end{itemize}
\end{frame}

% Worked Examples (5)
\begin{frame}{Worked Example 1}
\exampletitle{Roll a fair die. $P(\text{even})$?}\\[0.25em]
\worked
\begin{enumerate}
  \item Favorable outcomes: \{2,4,6\} (3 outcomes), total 6.
  \item $P(\text{even})=3/6=1/2$.
\end{enumerate}
\end{frame}

\begin{frame}{Worked Example 2}
\exampletitle{Draw one card from a standard 52-card deck. $P(\text{heart})$ and $P(\text{not heart})$?}\\[0.25em]
\worked
\begin{enumerate}
  \item Hearts: 13 of 52 $\Rightarrow P(\text{heart})=13/52=1/4$.
  \item Complement: $P(\text{not heart})=1-1/4=3/4$.
\end{enumerate}
\end{frame}

\begin{frame}{Worked Example 3}
\exampletitle{Two independent coin flips. $P(\text{at least one head})$?}\\[0.25em]
\worked
\begin{enumerate}
  \item Complement: $P(\text{no heads})=P(TT)=1/4$.
  \item So $P(\ge 1\text{ H})=1-1/4=3/4$.
\end{enumerate}
\end{frame}

\begin{frame}{Worked Example 4}
\exampletitle{Bag with 5 red and 3 blue. Draw 2 without replacement. $P(\text{both red})$?}\\[0.25em]
\worked
\begin{enumerate}
  \item First red: $5/8$. Second red given first red: $4/7$.
  \item $P= (5/8)(4/7)=20/56=5/14\approx0.357$.
\end{enumerate}
\end{frame}

\begin{frame}{Worked Example 5}
\exampletitle{In a class, $P(A)=0.4$, $P(B)=0.3$, $P(A\cap B)=0.1$. Find $P(A\cup B)$.}\\[0.25em]
\worked
\begin{enumerate}
  \item $P(A\cup B)=0.4+0.3-0.1=0.6$.
\end{enumerate}
\end{frame}

%%%%%%%%%%%%%%%%%%%%%%%%%%%%%%%%%%%%%%%%%%%%%%%%%%%%%%%%%%%%
\section{Lesson 6: Conditional Probability, Two-Way Tables \& Counting}
%%%%%%%%%%%%%%%%%%%%%%%%%%%%%%%%%%%%%%%%%%%%%%%%%%%%%%%%%%%%

\begin{frame}{Objectives}
\begin{itemize}
  \item Read and compute probabilities from two-way (contingency) tables.
  \item Compute conditional probability $P(A\mid B)$.
  \item Use permutations/combinations for counting.
\end{itemize}
\end{frame}

\begin{frame}{Key Ideas \& Formulas}
\begin{itemize}
  \item Two-way table organizes joint counts. Marginals on edges.
  \item Conditional probability: $P(A\mid B)=\dfrac{P(A\cap B)}{P(B)}$.
  \item Permutations (order matters): $_nP_r=\dfrac{n!}{(n-r)!}$.
  \item Combinations (order not): $_nC_r=\dfrac{n!}{r!(n-r)!}$.
\end{itemize}
\end{frame}

% Worked Examples (5)
\begin{frame}{Worked Example 1}
\exampletitle{Two-way table of 100 students:}\\[0.25em]
\small
\begin{center}
\begin{tabular}{l|ccc|c}
 & Like Math & Neutral & Dislike & Total\\\hline
Freshman & 18 & 7 & 5 & 30\\
Sophomore & 15 & 10 & 5 & 30\\
Upperclass & 20 & 12 & 8 & 40\\\hline
Total & 53 & 29 & 18 & 100\\
\end{tabular}
\end{center}
\worked
\begin{enumerate}
  \item $P(\text{Like})=53/100=0.53$.
  \item $P(\text{Freshman})=30/100=0.30$.
\end{enumerate}
\end{frame}

\begin{frame}{Worked Example 2}
\exampletitle{From the table, find $P(\text{Like}\mid \text{Freshman})$.}\\[0.25em]
\worked
\begin{enumerate}
  \item Within Freshmen, total $=30$, Like $=18$.
  \item $P(\text{Like}\mid \text{Freshman})=18/30=0.6$.
\end{enumerate}
\end{frame}

\begin{frame}{Worked Example 3}
\exampletitle{Are \textit{Freshman} and \textit{Like} independent?}\\[0.25em]
\worked
\begin{enumerate}
  \item Check if $P(\text{Like}\mid \text{Freshman})=P(\text{Like})$.
  \item $0.6 \neq 0.53$ $\Rightarrow$ not independent.
\end{enumerate}
\end{frame}

\begin{frame}{Worked Example 4}
\exampletitle{Passwords with letters A--Z (26) and digits 0--9 (10). How many 4-character codes with no repeats if (a) letters only; (b) letters or digits?}\\[0.25em]
\worked
\begin{enumerate}
  \item (a) Permutations without repeat: $26\cdot25\cdot24\cdot23={}_ {26}P_4$.
  \item (b) From 36 symbols: $36\cdot35\cdot34\cdot33={}_{36}P_4$.
\end{enumerate}
\end{frame}

\begin{frame}{Worked Example 5}
\exampletitle{A committee of 3 from 8 people. How many ways?}\\[0.25em]
\worked
\begin{enumerate}
  \item Order does not matter $\Rightarrow$ combinations: $_8C_3=\dfrac{8!}{3!\,5!}=\dfrac{8\cdot7\cdot6}{3\cdot2\cdot1}=56$.
\end{enumerate}
\end{frame}

%%%%%%%%%%%%%%%%%%%%%%%%%%%%%%%%%%%%%%%%%%%%%%%%%%%%%%%%%%%%
\section{Practice Sets \& Exit Tickets}
%%%%%%%%%%%%%%%%%%%%%%%%%%%%%%%%%%%%%%%%%%%%%%%%%%%%%%%%%%%%

\begin{frame}{Lesson 1 Practice (Assign)}
\begin{itemize}
  \item Identify type: (a) shoe size; (b) city of birth; (c) temperature.
  \item Make a relative-frequency table for: A,A,B,C,A,B,C,C.
  \item Propose an unbiased sampling plan to estimate average sleep hours.
\end{itemize}
\end{frame}

\begin{frame}{Lesson 2 Practice (Assign)}
\begin{itemize}
  \item Choose an appropriate display for survey data of favorite genres; justify.
  \item Given a histogram sketch, describe center, spread, shape, outliers.
\end{itemize}
\end{frame}

\begin{frame}{Lesson 3 Practice (Assign)}
\begin{itemize}
  \item Compute mean/median/mode/range for: 5,7,7,8,9,10.
  \item Find $Q_1,Q_3,$ IQR for: 2,3,4,6,7,8,12,14.
  \item Use IQR rule to test for outliers in: 1,2,2,3,4,80.
\end{itemize}
\end{frame}

\begin{frame}{Lesson 4 Practice (Assign)}
\begin{itemize}
  \item Draw a boxplot for: 12,14,15,16,18,21,22,24,30.
  \item Compare two boxplots with same median but different IQRs.
\end{itemize}
\end{frame}

\begin{frame}{Lesson 5 Practice (Assign)}
\begin{itemize}
  \item Deck of 52: $P(\text{face card})$; $P(\text{not diamond})$.
  \item Urn: 4 green, 6 yellow. Draw 2 without replacement: $P(\text{one of each})$.
\end{itemize}
\end{frame}

\begin{frame}{Lesson 6 Practice (Assign)}
\begin{itemize}
  \item From a two-way table, compute a conditional probability and decide independence.
  \item Locker combinations: digits 0--9, 3 digits, repetition allowed. How many?
\end{itemize}
\end{frame}

%%%%%%%%%%%%%%%%%%%%%%%%%%%%%%%%%%%%%%%%%%%%%%%%%%%%%%%%%%%%
\section{Summary \& Test Readiness}
%%%%%%%%%%%%%%%%%%%%%%%%%%%%%%%%%%%%%%%%%%%%%%%%%%%%%%%%%%%%

\begin{frame}{Unit Summary}
\begin{itemize}
  \item Choose displays that match data type; describe distributions.
  \item Mean/median/IQR tell center and spread; use IQR rule for outliers.
  \item Probability basics: complements, AND/OR, independence vs dependence.
  \item Conditional probability, two-way tables, and counting principles.
\end{itemize}
\end{frame}

\begin{frame}{Quick Exit Ticket}
\begin{enumerate}
  \item Data: 2,3,3,5,7,9. Find median and IQR.
  \item Bag: 3 red, 2 blue. Draw 2 without replacement. $P(\text{both blue})$?
  \item True/False: If $P(A\mid B)=P(A)$, then $A$ and $B$ are independent.
\end{enumerate}
\end{frame}

\begin{frame}
\centering \Large Questions?
\end{frame}

\end{document}
